\chapter{From Classical Normal Modes to Quantum Oscillators}

\begin{remarkbox}
This chapter follows the structure of the course plan in
\texttt{ideas\_Quantum\_Field\_Theory.md}. Its purpose is to turn the first
entry of the table of contents into a readable lecture: a free classical field
is decomposed into normal modes, each mode is recognized as a harmonic
oscillator, and the oscillator excitations become the first particle states of
quantum field theory.
\end{remarkbox}

\section{Starting Point: A Classical Field Decomposed into Modes}

Classical field theory describes a field by a function such as
\(\phi(t,\vect{x})\). For a free field, this function can be decomposed into
normal modes. This is the same idea that appears in mechanics when a system of
coupled oscillators is rewritten in terms of independent normal coordinates.

Begin with a finite chain of \(N\) identical masses connected by springs. If
\(q_j(t)\) is the displacement of the \(j\)-th mass, a typical quadratic
Lagrangian is
\[
    L = \sum_{j=1}^{N}\frac{m}{2}\dot q_j^2
        - \sum_{j=1}^{N}\frac{\kappa}{2}(q_{j+1}-q_j)^2.
\]
The variables \(q_j\) are coupled. After introducing normal coordinates
\(Q_r(t)\), one obtains the diagonal form
\[
    L = \sum_r
    \left(
        \frac{1}{2}\dot Q_r^{\,2}
        - \frac{1}{2}\omega_r^2 Q_r^2
    \right).
\]
The coupled chain is therefore equivalent to a collection of independent
oscillators.

A field is the continuum version of this situation. In a one-dimensional chain
with spacing \(a\), one may view
\[
    q_j(t) \approx \phi(t,x_j),
    \qquad x_j=ja.
\]
In the continuum limit, sums over sites become integrals over space, and the
normal-mode decomposition becomes a Fourier decomposition.

For a real scalar field in a box of volume \(V\), with periodic boundary
conditions, write schematically
\[
    \phi(t,\vect{x})
    = \frac{1}{\sqrt{V}}\sum_{\vect{k}}q_{\vect{k}}(t)
      e^{i\vect{k}\cdot\vect{x}},
\]
with the appropriate reality condition for a real field. The allowed momenta are
discrete because the field is placed in a finite box. This is a temporary
technical convenience; the continuum limit is obtained later by replacing sums
with integrals.

\begin{remarkbox}
The field is infinite-dimensional, but the free field is not mysterious at this
stage. The normal-mode decomposition rewrites it as a collection of independent
oscillating degrees of freedom.
\end{remarkbox}

\section{Each Mode as a Harmonic Oscillator}

For the free real scalar field, the classical Lagrangian is
\[
    L = \int \dd^3x\,
    \left(
        \frac{1}{2}\dot\phi^2
        - \frac{1}{2}(\nabla\phi)^2
        - \frac{1}{2}m^2\phi^2
    \right).
\]
Its equation of motion is the Klein--Gordon equation,
\[
    (\partial_t^2-\nabla^2+m^2)\phi(t,\vect{x})=0.
\]
After inserting the normal-mode expansion into the free Lagrangian and choosing
standard normalizations, the field decomposes into independent oscillator modes,
\[
    L = \sum_{\vect{k}}
    \left(
        \frac{1}{2}\dot q_{\vect{k}}^{\,2}
        - \frac{1}{2}\omega_{\vect{k}}^2 q_{\vect{k}}^{\,2}
    \right),
\]
where
\[
    \omega_{\vect{k}}=\sqrt{\vect{k}^{\,2}+m^2}
\]
in units with \(c=\hbar=1\).

The canonical momentum of the mode \(q_{\vect{k}}\) is
\[
    p_{\vect{k}}=\dot q_{\vect{k}},
\]
and the Hamiltonian becomes
\[
    H=\sum_{\vect{k}}H_{\vect{k}},
    \qquad
    H_{\vect{k}}
    =\frac{1}{2}p_{\vect{k}}^{\,2}
     +\frac{1}{2}\omega_{\vect{k}}^2q_{\vect{k}}^{\,2}.
\]
Thus each mode \(\vect{k}\) is a harmonic oscillator of frequency
\(\omega_{\vect{k}}\).

\begin{examplebox}
The classical field is not quantized all at once. First, the free classical
field is diagonalized into modes. Then each mode is quantized using the same
oscillator algebra known from elementary quantum mechanics.
\end{examplebox}

\section{Quantizing a Single Oscillator}

For one harmonic oscillator, the classical variables \(q\) and \(p\) become
operators \(\hat q\) and \(\hat p\) satisfying
\[
    [\hat q,\hat p]=i.
\]
The Hamiltonian is
\[
    \hat H
    =\frac{1}{2}\hat p^2
     +\frac{1}{2}\omega^2\hat q^2.
\]
Introduce the annihilation and creation operators
\[
    \hat a
    =\sqrt{\frac{\omega}{2}}\,\hat q
     +\frac{i}{\sqrt{2\omega}}\,\hat p,
    \qquad
    \hat a^\dagger
    =\sqrt{\frac{\omega}{2}}\,\hat q
     -\frac{i}{\sqrt{2\omega}}\,\hat p.
\]
Equivalently,
\[
    \hat q=\frac{1}{\sqrt{2\omega}}(\hat a+\hat a^\dagger),
    \qquad
    \hat p=-i\sqrt{\frac{\omega}{2}}(\hat a-\hat a^\dagger).
\]
The canonical commutator implies
\[
    [\hat a,\hat a^\dagger]=1.
\]
The Hamiltonian becomes
\[
    \hat H=\omega\left(\hat a^\dagger\hat a+\frac{1}{2}\right).
\]
This is the algebraic form that will be copied for every field mode.

\section{Creation and Annihilation Operators}

The operator \(\hat a^\dagger\) raises the excitation number of the oscillator,
while \(\hat a\) lowers it. If \(|0\rangle\) is the oscillator ground state,
then
\[
    \hat a|0\rangle=0.
\]
Repeated action of \(\hat a^\dagger\) produces excited states,
\[
    |n\rangle=\frac{(\hat a^\dagger)^n}{\sqrt{n!}}|0\rangle,
    \qquad n=0,1,2,\ldots.
\]
In a single oscillator problem, this is a convenient way to organize the energy
spectrum. In a field theory, the same algebra acquires a new interpretation:
creation operators create field quanta.

\section{Number Operator and Energy Spectrum}

The number operator is
\[
    \hat N=\hat a^\dagger\hat a.
\]
It satisfies
\[
    \hat N|n\rangle=n|n\rangle.
\]
The Hamiltonian eigenvalue equation is
\[
    \hat H|n\rangle
    =\omega\left(n+\frac{1}{2}\right)|n\rangle.
\]
The integer \(n\) counts oscillator excitations. The term
\(\omega/2\) is the zero-point energy.

For a field, there will be one number operator for each mode,
\[
    \hat N_{\vect{k}}
    =\hat a_{\vect{k}}^\dagger\hat a_{\vect{k}}.
\]
The eigenvalue of \(\hat N_{\vect{k}}\) counts how many quanta occupy the mode
\(\vect{k}\).

\section{The Vacuum State of One Oscillator}

The oscillator vacuum is defined algebraically by
\[
    \hat a|0\rangle=0.
\]
It is not the state of zero energy; its energy is
\[
    E_0=\frac{1}{2}\omega.
\]
This small fact becomes conceptually important in field theory. A free field has
one oscillator for every allowed momentum mode, so the formal vacuum energy is
\[
    E_{\mathrm{vac}}
    =\frac{1}{2}\sum_{\vect{k}}\omega_{\vect{k}}.
\]
In the continuum limit this expression diverges. Later chapters will explain
normal ordering, regularization, renormalization, and the physical meaning of
vacuum energy. At this stage the important lesson is simpler: the quantum vacuum
is a state with no excitations, not a state with no structure.

\section{From One Oscillator to Infinitely Many Oscillators}

For the scalar field, each normal mode receives its own creation and
annihilation operators,
\[
    \hat a_{\vect{k}},
    \qquad
    \hat a_{\vect{k}}^\dagger.
\]
They satisfy
\[
    [\hat a_{\vect{k}},\hat a_{\vect{k}'}^\dagger]
    =\delta_{\vect{k},\vect{k}'},
    \qquad
    [\hat a_{\vect{k}},\hat a_{\vect{k}'}]=0,
    \qquad
    [\hat a_{\vect{k}}^\dagger,\hat a_{\vect{k}'}^\dagger]=0.
\]
The free-field vacuum is the state annihilated by every annihilation operator,
\[
    \hat a_{\vect{k}}|0\rangle=0
    \qquad \text{for every }\quad \vect{k}.
\]
A one-particle state of momentum \(\vect{k}\) is
\[
    |\vect{k}\rangle
    =\hat a_{\vect{k}}^\dagger|0\rangle.
\]
A two-particle state may be written as
\[
    |\vect{k}_1,\vect{k}_2\rangle
    =\hat a_{\vect{k}_1}^\dagger
     \hat a_{\vect{k}_2}^\dagger|0\rangle,
\]
up to normalization conventions.

\begin{definitionbox}
The \emph{Fock space} of a free quantum field is the Hilbert space generated
from the vacuum by applying creation operators for all allowed modes. It is the
natural state space for a theory in which the number of particles is not fixed.
\end{definitionbox}

\section{Field Quanta as Mode Excitations}

In the free theory, the Hamiltonian can be written as
\[
    \hat H
    =\sum_{\vect{k}}\omega_{\vect{k}}
    \left(
        \hat a_{\vect{k}}^\dagger\hat a_{\vect{k}}
        +\frac{1}{2}
    \right).
\]
The first term counts excitations of each mode. The state
\(\hat a_{\vect{k}}^\dagger|0\rangle\) has energy \(\omega_{\vect{k}}\) above
the free vacuum and momentum \(\vect{k}\). It is interpreted as a one-particle
state.

The field operator itself is reconstructed from the same mode operators. In a
standard continuum normalization,
\[
    \hat\phi(t,\vect{x})
    =\int \frac{\dd^3k}{(2\pi)^3}
      \frac{1}{\sqrt{2\omega_{\vect{k}}}}
      \left(
          \hat a_{\vect{k}}e^{-i\omega_{\vect{k}}t+i\vect{k}\cdot\vect{x}}
          +\hat a_{\vect{k}}^\dagger
           e^{i\omega_{\vect{k}}t-i\vect{k}\cdot\vect{x}}
      \right).
\]
The coefficients of the classical mode expansion have become operators. The
field operator can annihilate and create quanta because it contains both
\(\hat a_{\vect{k}}\) and \(\hat a_{\vect{k}}^\dagger\).

\begin{remarkbox}
A field quantum is not a tiny localized lump of classical field. In the free
theory, it is a mode excitation created by a creation operator. Localization,
wave packets, interactions, and particle detection require additional steps.
\end{remarkbox}

\section{What Is New Compared with Classical Field Theory?}

Classical field theory assigns numbers to field configurations. Quantum field
theory assigns operators to fields and vectors to states. The basic object is no
longer a definite field profile \(\phi(x)\), but an operator \(\hat\phi(x)\)
acting on a Hilbert space.

The free quantum field still satisfies the free equation of motion,
\[
    (\Box+m^2)\hat\phi(x)=0,
\]
but the meaning has changed. This is now an equation for an operator-valued
distribution, not for a classical function.

The conceptual dictionary of the chapter is:
\[
\begin{array}{ccl}
\text{classical normal mode} & \longrightarrow & \text{quantum oscillator},\\[0.25em]
\text{oscillator excitation} & \longrightarrow & \text{field quantum},\\[0.25em]
\text{many oscillator excitations} & \longrightarrow & \text{Fock-space state},\\[0.25em]
\text{classical field} & \longrightarrow & \text{field operator}.
\end{array}
\]

\begin{summarybox}
A free classical field can be decomposed into independent normal modes. Each
mode is a harmonic oscillator. Quantizing all of these oscillators produces
creation and annihilation operators, a vacuum state, number operators, and Fock
space. In the free theory, particles are mode excitations of quantum fields.
This is the starting point for propagators, perturbation theory, and scattering.
\end{summarybox}
